\chapter{Tables of Reasoning Methods and Evidence}
\label{sec:p14-reasoning-methods-and-evidence}
\index{reasoning methods!comparison}\index{evidence!interpretation}

The tables in this chapter compare the question answered, the object
returned and the conditions on which that answer depends. A truth-preserving
derivation, an acceptable argument, a probability assignment and a decision
can concern the same proposition while establishing different things.
The linked entries supply the exact definitions and results behind each row.
The task finder in \cref{sec:p14-reasoning-task-finder} locates an entry from
a question; the present chapter helps interpret and compare its output.

\section{Deductive validity, defeasible acceptability, and inductive support}
\label{sec:p14-evidence-support}

For deduction, name the language, calculus and interpretations. For
defeasibility, additionally name the preference, default or attack mechanism
and the policy for multiple extensions. For statistical support, name the
probability model, observations and comparison. These specifications determine
what counts as a counterexample or a defeating consideration
\cite{Enderton2001,KrausLehmannMagidor1990,Dung1995}.

\begin{longtable}{@{}p{.18\textwidth}p{.29\textwidth}p{.43\textwidth}@{}}
\caption{Deductive validity, defeasible acceptability, and inductive support}\\
\toprule
Relation or method & What the output establishes & Conditions and boundary \\
\midrule
\endfirsthead
\caption[]{Deductive validity, defeasible acceptability, and inductive support (continued)}\\
\toprule
Relation or method & What the output establishes & Conditions and boundary \\
\midrule
\endhead
\midrule\multicolumn{3}{r}{\emph{continued on next page}}\\
\endfoot
\endlastfoot
Classical semantic consequence & Every interpretation satisfying the
premises satisfies the conclusion. & Fix the interpretation class and
truth clauses. An inconsistent classical premise set entails every formula;
validity of an argument does not establish premise truth.
\cref{sec:prop-syntax-semantics-02,sec:fol-semantics-02}. \\
Derivation in a calculus & A finite derivation uses the licensed rules and
open assumptions. & Soundness connects this syntactic fact to the specified
semantics; completeness supplies the converse for its stated inputs.
\cref{thm:fol-soundness,cor:fol-completeness-countable} and the exact
\hyperref[sec:p14-rule-sheets]{calculus sheets}. \\
Preferential consequence & The conclusion holds in the minimal premise
states of the supplied model. & Smoothness and ranking impose different
conditions; added information can change which states are minimal.
\cref{def:p08-preferential-model,def:p08-nonmonotonic}. \\
Argument acceptability & Membership in the selected extension, or membership
in some or every extension under the named semantics. & The input attack
graph supplies the conflicts. Acceptance does not establish factual truth;
stable extensions can be absent.
\cref{def:p18-finite-framework,def:p18-extension-semantics}. \\
Logical abduction & Permitted assumptions are jointly satisfiable with the
background and entail the observations. & Positive atomic hypotheses and
observations are the finite input here. Subset-minimality and minimum
cardinality differ. Neither establishes the hypotheses from the observations.
\cref{def:p16-abduction-models,thm:p16-abduction-truth-boundary}. \\
Bayesian support & A posterior, or a posterior-odds change, follows from the
declared joint model after conditioning. & Proper normalization and the
observation mechanism matter. A supported hypothesis need not be entailed,
and a prior-zero hypothesis cannot gain positive mass on the ordinary
conditioning domain.
\cref{thm:p15-bayes-posterior-kernel,thm:p15-model-posterior-odds}. \\
Empirical support for an account of reasoning & Recorded responses constrain
a claim about the specified participants, task or relation. & Distinguish
score, construct, estimate, design and normative benchmark. A correct answer
does not uniquely identify its cognitive process.
\cref{def:p19-empirical-question,def:p19-evidence-record}. \\
\bottomrule
\end{longtable}

Classical entailment is monotone in its premises: adding premises restricts
the satisfying interpretations. Preferential consequence can withdraw a
conclusion because adding information changes the selected minimal states.
This is a difference between the defined relations. It does not make every
nonmonotonic inference an estimate of probability, or make a probability
update a preferential inference. The classical inconsistent-input case is
worked out for explanatory entailment in
\cref{thm:p16-abduction-semantics}; nonexplosive alternatives have their
own designation or derivation conditions in
\cref{def:p08-paraconsistency,def:p08-lp-nonexplosion}.

\begin{longtable}{@{}p{.24\textwidth}p{.32\textwidth}p{.34\textwidth}@{}}
\caption{Witnesses of failure and what a successful check leaves open}\\
\toprule
Assessment & Witness of failure & What a successful check leaves open \\
\midrule
\endfirsthead
\caption[]{Witnesses of failure and what a successful check leaves open (continued)}\\
\toprule
Assessment & Witness of failure & What a successful check leaves open \\
\midrule
\endhead
\midrule\multicolumn{3}{r}{\emph{continued on next page}}\\
\endfoot
\endlastfoot
Classical consequence & A satisfying interpretation of all premises that
falsifies the conclusion. & Whether the premises describe the intended
situation, and whether the formalization preserves its meaning. \\
An offered derivation & A specific unlicensed inference, unmet side
condition or incorrect dependency. & Another derivation can still exist;
rejecting this proof does not refute its conclusion. \\
Abductive explanation & Either inconsistency, a countervaluation to an
observation, or an ineligible hypothesis. & Another permitted explanation
can exist; a minimal explanation need not be true. \\
Stable-extension membership & An extension omitting the argument refutes
skeptical membership when at least one extension exists. & If no stable
extension exists, some-extension membership fails and every-extension
membership is vacuous unless an existence requirement is imposed. \\
A model's selected predictive feature & A specified discrepancy that is
atypical of the selected replications. & Other aspects may fit or fail;
the check does not supply posterior model odds. \\
\bottomrule
\end{longtable}

The last two rows require particular care about quantifiers. The existence
and separation examples for finite argumentation are in
\cref{sec:p18-finite-argumentation}; the replication and discrepancy
definitions for model criticism are
\cref{def:p15-ppc-model,def:p15-ppc-discrepancy}. A failed search without
the guarantee needed to exhaust its target class supplies no such universal
negative conclusion.

\section{Probability, possibility, belief, and other uncertainty measures}
\label{sec:p14-evidence-uncertainty}

The following comparison uses a finite outcome space and distinguishes
event weights from qualitative modal belief. Its probability foundation is
\cref{def:p15-finite-probability-space}. The alternative event-weight
definitions follow Halpern's finite presentations in Chapter 2 of
\emph{Reasoning about Uncertainty}
\cite[secs.~2.3--2.5]{HalpernUncertaintyChapter2Extract}.

\begin{hbdefinition}{def:p14-evidence-uncertainty-data}{Four finite event-weight constructions}
Let \(W\) be finite and nonempty, and let \(A\subseteq W\).
Write \(A^c=W\setminus A\). A probability mass \(p\) has
\(p(w)\geq0\) and \(\sum_{w\in W}p(w)=1\), with
\(P_p(A)=\sum_{w\in A}p(w)\).

For a nonempty family \(\mathcal C\) of such masses, define the envelopes
\[
 \underline P_{\mathcal C}(A)=\inf_{p\in\mathcal C}P_p(A),\qquad
 \overline P_{\mathcal C}(A)=\sup_{p\in\mathcal C}P_p(A).
\]
For \(\pi:W\to[0,1]\) with \(\max_{w\in W}\pi(w)=1\), define
\[
 \operatorname{Poss}_{\pi}(A)=\max_{w\in A}\pi(w),\qquad
 \operatorname{Nec}_{\pi}(A)=1-\operatorname{Poss}_{\pi}(A^c),
\]
where the empty maximum is zero.
For \(m:2^W\to[0,1]\) with \(m(\varnothing)=0\) and
\(\sum_{B\subseteq W}m(B)=1\), define
\[
 \operatorname{Bel}_m(A)=\sum_{B\subseteq A}m(B),\qquad
 \operatorname{Pl}_m(A)=\sum_{B:B\cap A\ne\varnothing}m(B).
\]
These are respectively the mass-based belief and plausibility functions.
An allocation for \(m\) chooses a probability mass \(q_B\) on each
nonempty \(B\) with \(m(B)>0\), extended by zero off \(B\).
Let \(\mathcal C_m\) contain every mass
\(p(w)=\sum_{B:m(B)>0}m(B)q_B(w)\) obtained this way.
\end{hbdefinition}

\begin{hbtheorem}{thm:p14-evidence-finite-measures}{proved=annex:p14.evidence.finite-measures}{Finite duality, union laws and allocated bounds}
All four constructions are normalized and monotone under event inclusion.
For \(A,B\subseteq W\):
\begin{enumerate}
\item \(\overline P_{\mathcal C}(A)=1-\underline P_{\mathcal C}(A^c)\).
      For disjoint \(A,B\), the lower envelope is superadditive and the
      upper envelope is subadditive.
\item \(\operatorname{Poss}_{\pi}(A\cup B)=
      \max\{\operatorname{Poss}_{\pi}(A),\operatorname{Poss}_{\pi}(B)\}\)
      and \(\operatorname{Nec}_{\pi}(A\cap B)=
      \min\{\operatorname{Nec}_{\pi}(A),\operatorname{Nec}_{\pi}(B)\}\).
\item \(\operatorname{Pl}_m(A)=1-\operatorname{Bel}_m(A^c)\), and
      \(\operatorname{Bel}_m(A)\leq\operatorname{Pl}_m(A)\).
      Their difference is exactly
      \(\sum_{B:B\cap A\ne\varnothing,\ B\cap A^c\ne\varnothing}m(B)\).
\item \(\mathcal C_m\) is nonempty, and its event envelopes satisfy
      \(\underline P_{\mathcal C_m}(A)=\operatorname{Bel}_m(A)\) and
      \(\overline P_{\mathcal C_m}(A)=\operatorname{Pl}_m(A)\).
      Both bounds are attained, possibly at different allocations.
\item \(\operatorname{Bel}_m(A\cup B)\geq
      \operatorname{Bel}_m(A)+\operatorname{Bel}_m(B)
      -\operatorname{Bel}_m(A\cap B)\).
      If positive mass is confined to singleton subsets, belief and
      plausibility both equal the corresponding probability measure.
\end{enumerate}
No updating, evidence-combination or decision rule is part of these claims.
\end{hbtheorem}

\begin{longtable}{@{}p{.22\textwidth}p{.31\textwidth}p{.37\textwidth}@{}}
\caption{Probability, possibility, belief, and other uncertainty measures}\\
\toprule
Representation & Input and event calculation & Interpretive boundary \\
\midrule
\endfirsthead
\caption[]{Probability, possibility, belief, and other uncertainty measures (continued)}\\
\toprule
Representation & Input and event calculation & Interpretive boundary \\
\midrule
\endhead
\midrule\multicolumn{3}{r}{\emph{continued on next page}}\\
\endfoot
\endlastfoot
One probability law & Add nonnegative outcome masses totaling one. &
Probabilities of disjoint events add. State whether the law describes a
sampling mechanism, a joint model or a subjective assessment. \\
Set of probability laws & Take the event's infimum and supremum over the
specified nonempty family. & An interval is a range across retained laws;
it is not a posterior distribution over that family. \\
Finite possibility & Take the maximum of the supplied singleton scores. &
Scores normalize by a maximum; summing them need not give one. \\
Finite belief and plausibility & Aggregate supplied subset masses by
containment or nonempty intersection. & Unallocated mass can give a gap
between the two values. It does not specify an additive split within a subset. \\
Qualitative modal belief & Require truth at the accessible worlds under the
selected doxastic semantics. & It has no numerical event weight merely by
being called belief. The \(\mathsf{KD45}_{I}\) card does not assert factivity.
\cref{entry:p08-kd45n}. \\
\bottomrule
\end{longtable}

\begin{hbexample}{def:p14-evidence-measure-separation}{Exact weights on the same eight events}
Take \(W=\{a,b,c\}\), \(p=(1/2,1/5,3/10)\), and
\(\mathcal C=\{(t,7/10-t,3/10):0\leq t\leq7/10\}\).
Set \(m(\{a,b\})=7/10\), \(m(\{c\})=3/10\), with all other
masses zero; then \(\mathcal C=\mathcal C_m\). Independently set
\(\pi(a)=1\), \(\pi(b)=3/5\), \(\pi(c)=1/5\).
These are specified mathematical inputs, not estimates from observations.
\end{hbexample}

\begin{hbtheorem}{thm:p14-evidence-measure-separation}{proved=annex:p14.evidence.measure-separation}{The complete finite comparison and a non-belief lower envelope}
For the preceding inputs the values are those in
\cref{tab:p14-evidence-eight-events}. In particular the probability,
possibility and necessity values for \(\{a,b\}\) are respectively
\(7/10,1,4/5\); no one of these replaces either of the others.

Moreover, a general lower probability need not be a mass-based belief
function. On \(\{1,2,3,4\}\), take the two retained masses
\((2/5,0,0,3/5)\) and \((0,1/2,1/2,0)\).
For \(A=\{1,2\}\), \(B=\{1,3\}\), their lower envelope gives
\[
 \underline P(A)=\underline P(B)=\underline P(A\cup B)=2/5,
 \qquad \underline P(A\cap B)=0.
\]
It violates the inequality in item 5 of
\cref{thm:p14-evidence-finite-measures} and hence cannot equal
\(\operatorname{Bel}_m\) for any mass assignment of that form.
\end{hbtheorem}

\begin{table}[htbp]
\centering\small
\begin{tabular}{@{}lccccc@{}}
\toprule
Event & \(P_p\) & \(\operatorname{Bel}_m\) &
\(\operatorname{Pl}_m\) &
\(\operatorname{Poss}_{\pi}\) & \(\operatorname{Nec}_{\pi}\) \\
\midrule
\(\varnothing\) & 0 & 0 & 0 & 0 & 0 \\
\(\{a\}\) & \(1/2\) & 0 & \(7/10\) & 1 & \(2/5\) \\
\(\{b\}\) & \(1/5\) & 0 & \(7/10\) & \(3/5\) & 0 \\
\(\{c\}\) & \(3/10\) & \(3/10\) & \(3/10\) & \(1/5\) & 0 \\
\(\{a,b\}\) & \(7/10\) & \(7/10\) & \(7/10\) & 1 & \(4/5\) \\
\(\{a,c\}\) & \(4/5\) & \(3/10\) & 1 & 1 & \(2/5\) \\
\(\{b,c\}\) & \(1/2\) & \(3/10\) & 1 & \(3/5\) & 0 \\
\(W\) & 1 & 1 & 1 & 1 & 1 \\
\bottomrule
\end{tabular}
\caption{Different event-weight constructions under the stated inputs.
The additive law belongs to the displayed family. Its lower and upper
envelopes equal the belief and plausibility columns. The possibility scores
are a separate stipulated representation.}
\label{tab:p14-evidence-eight-events}
\end{table}

An event interval alone does not choose an action. A decision additionally
needs consequences, preferences or loss, feasible actions, and a rule for
using the retained laws. The particular finite maxmin and robust-dominance
rules in \cref{def:p17-finite-maxmin,def:p17-robust-dominance} give different
comparisons even when they use the same prior list. Fuzzy truth values in
\cref{def:p08-tnorm-residuum} answer a further distinct semantic
question; they do not become probabilities by taking values in \([0,1]\).

\section{Estimation, testing, prediction, and model criticism}
\label{sec:p14-evidence-statistical}

Fix the observed variable, the target quantity and the probability law in
each claim. The posterior concerns an unknown quantity conditional on the
observations. A sampling guarantee concerns the random procedure at a fixed
parameter. Prediction concerns a future observation under its specified law.
A check compares an observed feature with its replication distribution
\cite{Neyman1937Confidence,KassRaftery1995BayesFactors,GelmanMengStern1996Main}.

\begin{longtable}{@{}p{.18\textwidth}p{.32\textwidth}p{.40\textwidth}@{}}
\caption{Estimation, testing, prediction, and model criticism}\\
\toprule
Quantity & Exact question or defining calculation & Conditions and what it does not say \\
\midrule
\endfirsthead
\caption[]{Estimation, testing, prediction, and model criticism (continued)}\\
\toprule
Quantity & Exact question or defining calculation & Conditions and what it does not say \\
\midrule
\endhead
\midrule\multicolumn{3}{r}{\emph{continued on next page}}\\
\endfoot
\endlastfoot
Estimate & A declared function of the observations reports a value for a
specified estimand. & The estimand fixes population, units and comparison.
Point estimation by itself gives neither an error bound nor causal
identification. \cref{def:p19-population-and-estimand}. \\
Likelihood & Evaluate a sampling density at the observed data as a function
of its parameter. & It need not integrate to one over parameters and is not
a posterior. Fix compatible versions and observation descriptions.
\cref{def:p15-bayes-experiment,def:p15-model-comparison-experiment}. \\
Simple likelihood-ratio test & Compare two fixed densities by the stated
threshold and boundary randomization. & Most-powerful level control concerns
the specified simple pair; a uniformly most-powerful test for a composite
alternative needs its own argument.
\cref{thm:p15-neyman-pearson,thm:p15-binomial-most-powerful}. \\
Bayes factor & Ratio of the two marginal data densities, each integrating
its own proper parameter prior. & Retains within-model priors. Maximizing
instead of integrating changes the calculation. Zero denominators and
improper-prior constants need separate treatment.
\cref{thm:p15-model-posterior-odds,thm:p15-improper-model-priors}. \\
Posterior model probability & Weight a marginal data density by its model
prior and normalize across models. & Depends on model prior probabilities
as well as Bayes factors. It concerns the supplied alternatives, not a
certificate that one of them describes the world exactly.
\cref{def:p15-model-comparison-experiment}. \\
Valid \(p\)-value & Under each declared null law, probability of obtaining a
value at most \(t\) is at most \(t\). & Tail choice, ties, observation
mechanism and the null family matter. The number is not a posterior
probability of the null. \cref{def:p15-valid-p-value,thm:p15-discrete-p-values}. \\
Confidence set & The procedure includes the fixed target with at least the
stated probability for every parameter in its scope. & Probability ranges
over repeated data and declared randomization. Interval shape and exact
coverage require additional properties; discrete coverage can exceed the
nominal level.
\cref{def:p15-confidence-procedure,thm:p15-binomial-confidence-coverage}. \\
Credible set & The posterior assigns the reported set its stated conditional
content. & A fixed-parameter sampling guarantee does not follow from that
content. Proper joint models yield the separate prior-average calibration
identity. \cref{def:p15-bayes-credible-set,thm:p15-bayes-coverage}. \\
Predictive distribution & Integrate a specified future-observation kernel
against the posterior. & Include parameter uncertainty and state dependence
on the earlier observations. Its probability is conditional on the declared
model. \cref{thm:p15-bayes-predictive-kernel}. \\
Posterior predictive check & Compare discrepancies of observed and replicated
data, using the selected shared parameter and replication scheme. & It tests
the chosen aspect. Its reference probability need not be a valid classical
\(p\)-value under a fixed null sampling law. Pivot and prior-model
calibration have separate conditions.
\cref{thm:p15-ppc-joint,thm:p15-ppc-pivot,thm:p15-ppc-normal-mean}. \\
Measurement uncertainty interval & Multiply a combined standard uncertainty
by a stated coverage factor and report the interval. & A multiplier does
not determine a universal coverage probability. Include shared influences
and the distributional interpretation.
\cref{def:p15-measurement-expanded,thm:p15-measurement-coverage-models}. \\
\bottomrule
\end{longtable}

\subsection{Three comparisons that change the interpretation}

First, \cref{thm:p15-model-posterior-odds} separates a factor from odds.
A numerator-model factor of 20, with prior odds \(1:100\), gives posterior
odds \(1:5\) and posterior probability \(1/6\). The same factor with
equal prior odds gives probability \(20/21\). Stating the factor alone
leaves the prior odds unspecified. Its logarithm also compares marginal
predictive densities, whereas a loss function determines whether either
posterior comparison warrants an action.

Second, the exact Bernoulli calculation in \cref{thm:p15-bayes-coverage}
gives posterior intervals of content \(0.95\) after either outcome, yet
their repeated-sampling coverage at the fixed parameter \(0.99\) is zero.
This is compatible with the proved prior-average coverage identity. The
uniform guarantee of a confidence procedure and the content of a credible
set have different quantifiers. Identical-looking interval endpoints do
not identify the two claims; each interpretation needs its own calculation.

Third, the normal location check in \cref{thm:p15-ppc-normal-mean} has
reference probability \(1/2\) for every observed sample under its stated
model and prior convention. Thus absence of an extreme reference value
can arise because the discrepancy does not constrain that fitted feature.
Changing the discrepancy or replication asks a different question. Compare
the conditional dispersion calculation in \cref{thm:p15-ppc-poisson}; a
correct fitted total can coexist with failure in the counts' dispersion.

\subsection{Training, evaluation and repeated inspection}

The finite-class risk bound in \cref{thm:p15-conc-finite-learning} assumes
the displayed independent identically distributed training observations,
bounded loss and a fixed finite candidate class. Its simultaneous control
permits the stated selection within that class. A bound for one fixed
candidate cannot be attached unchanged to whichever candidate looks best
after searching many alternatives. The explicit selection failure is
\cref{thm:p15-conc-selection-example}; independent validation has the
different guarantee in \cref{thm:p15-conc-validation}.
Likewise repeated inspection of a fixed-sample interval does not create a
time-uniform procedure. The simultaneous finite-range construction in
\cref{thm:p15-conc-time-uniform} supplies such a guarantee by its own
allocation of error across sample sizes. Model misspecification, changing
populations and unrecorded selection remain questions about whether the
displayed assumptions hold.

\section{Causal identification, explanatory comparison, and decision criteria}
\label{sec:p14-evidence-causal-explanatory-decision}

An explanatory hypothesis, a causal effect and a preferred action require
different structures. Logical abduction tests assumptions against a theory
and observations. Intervention semantics describes how specified mechanisms
change. A decision representation combines consequences and probabilities
with preferences or loss. These operations do not supply one another's
missing inputs \cite{EiterGottlobAbductionManuscript,Pearl2009CausalOverview,
Rubinstein2015LectureNotes}.

\begin{longtable}{@{}p{.20\textwidth}p{.31\textwidth}p{.39\textwidth}@{}}
\caption{Causal identification, explanatory comparison, and decision criteria}\\
\toprule
Criterion & Required structures and conclusion & Comparison or failure boundary \\
\midrule
\endfirsthead
\caption[]{Causal identification, explanatory comparison, and decision criteria (continued)}\\
\toprule
Criterion & Required structures and conclusion & Comparison or failure boundary \\
\midrule
\endhead
\midrule\multicolumn{3}{r}{\emph{continued on next page}}\\
\endfoot
\endlastfoot
Abductive minimality & Find a permitted, consistent, observationally
sufficient hypothesis set; minimize inclusion or cardinality. & Different
incomparable minimal sets can remain. Fewer assumptions alone supplies no
probability or causal ordering.
\cref{thm:p16-abduction-minimality,thm:p16-abduction-examples}. \\
Deductive-nomological explanation & Laws and conditions entail the result;
the source additionally requires the indicated empirical content, role of
laws and truth conditions. & Bare entailment is insufficient for the
1948 account. Its explanatory use does not identify a causal direction.
\cref{def:p16-dn-schema,sec:p16-deductive-nomological-explanation}. \\
Conditioning & Retain the observational mechanisms and reweight outcomes
consistent with the conditioning information. & Positive-probability event
conditioning and continuous conditional kernels have their distinct domains.
Conditioning on a treatment is generally different from intervening on it.
\cref{def:p15-finite-conditioning,def:p16-causal-intervention}. \\
Intervention calculation & Replace the selected structural equations,
preserving the unreplaced mechanisms and exogenous law as stipulated. &
Computes the supplied model's effect. A fitted observational table alone
does not supply the equations or their invariance.
\cref{thm:p16-causal-intervention-laws,def:p16-causal-finite-calculation}. \\
Causal identification & Every model in the declared class with the same
observational law has the same target intervention law. & State the graph,
latent structure, support and alphabet quantifiers. Estimation error is a
further question; a nonidentified effect can remain so with unlimited
observational data. \cref{thm:p16-causal-nonidentification,
def:p16-graph-uniform-positive}. \\
Expected utility & A supplied lottery preference is represented by the
expected value of consequence utility. & Finite-prize completeness,
transitivity, mixture independence and continuity characterize the displayed
representation. Ordinal ordering of sure prizes alone does not fix it.
\cref{thm:p17-finite-vnm,thm:p17-vnm-affine-uniqueness}. \\
Posterior-risk choice & Among feasible actions, minimize expected loss
under the posterior. & Loss, feasible actions and timing are inputs.
Optimality is conditional on them, and minimizers require the theorem's
finite-action or attainment assumptions.
\cref{def:p15-bayes-loss,thm:p15-bayes-decision-rule}. \\
Maxmin expected utility & Maximize the worst expected utility over the
specified retained priors. & Worst expected utility is different from
worst utility over raw states. Enlarging the prior set can change rankings.
\cref{def:p17-finite-maxmin,thm:p17-finite-maxmin-properties}. \\
Robust dominance & One act's expected utility is at least the other's
under every retained prior. & This comparison can be incomplete; a maxmin
ranking need not imply robust dominance.
\cref{thm:p17-robust-dominance}. \\
Nash equilibrium & No player benefits by a unilateral deviation while the
others' mixed strategies stay fixed. & Cardinal payoffs and the specified
game are inputs. The condition implies neither collective optimality nor
uniqueness or convergence.
\cref{def:p17-nash-equilibrium,thm:p17-best-response-geometry}. \\
\bottomrule
\end{longtable}

An identification formula is a conclusion under a model class, not evidence
that the class fits a study. Back-door adjustment requires its exact graphical
and support conditions in \cref{sec:p16-graphical-separation-adjustment}.
The graph-uniform criterion in \cref{thm:p16-graph-uniform-equivalences}
quantifies over every permitted finite alphabet family. Its failure witness
does not automatically prove failure for a separately restricted binary
family. Conversely, algebraic evaluation in one known structural model
does not prove identification across other compatible models.

A recommendation adds a further judgment about consequences. Changing
utility or loss can change the action while leaving the probability model
unchanged. Positive affine changes preserve the expected-utility comparisons
of \cref{thm:p17-vnm-affine-uniqueness}; arbitrary increasing changes of
individual consequence utilities need not. Empirical frequencies of choices
do not decide which preferences or losses ought to govern a recommendation.
The task, measure and benchmark distinction in \cref{def:p19-benchmark}
applies equally when a study labels responses rational or irrational.

\section{Assumptions, guarantees, failure cases, and method alternatives}
\label{sec:p14-evidence-assumptions-alternatives}

An alternative is useful when it addresses the same target under defensible
changed assumptions, or when it makes a different target explicit. The
following cases identify the missing dependency and what the proposed next
method would need. They give no universal ordering of methods.

\begin{longtable}{@{}p{.23\textwidth}p{.31\textwidth}p{.36\textwidth}@{}}
\caption{Assumptions, guarantees, failure cases, and method alternatives}\\
\toprule
Situation & Guarantee available or lost & Alternative and its obligation \\
\midrule
\endfirsthead
\caption[]{Assumptions, guarantees, failure cases, and method alternatives (continued)}\\
\toprule
Situation & Guarantee available or lost & Alternative and its obligation \\
\midrule
\endhead
\midrule\multicolumn{3}{r}{\emph{continued on next page}}\\
\endfoot
\endlastfoot
Classical premises are inconsistent & Ordinary entailment becomes vacuous;
abductive consistency fails. & Check premises and their formalization, or
use a declared nonexplosive calculus. Its consistency or designation
mechanism must be specified.
\cref{thm:p16-abduction-semantics,sec:p08-paraconsistent-and-nonmonotonic}. \\
An exception defeats a normal-case conclusion & The relevant preferential
selection or extension can change. & Revise the premises or priorities and
recompute the selected semantics. Adding an exception must be justified;
asking about one does not establish it.
\cref{def:p08-preferential-model,def:p18-scheme-assessment-terms}. \\
Additional records duplicate earlier information & Multiplying independent
likelihood contributions misstates the joint observation model. & Specify
conditional likelihoods. The duplicated-flag calculation and sequential
update identity are
\cref{thm:p15-bayes-flag-calculation,thm:p15-bayes-sequential-update}. \\
The observation was selected & An ordinary unselected likelihood can give
the wrong posterior. & Include the mechanism or justify its cancellation
under the exact parameter-independence condition.
\cref{thm:p15-bayes-stopping}. \\
No defensible single prior has been supplied & An ordinary posterior or
expected-utility answer is unspecified. & State a sensitivity family or
retained prior set, and specify its update or decision rule. Event envelopes
alone do not choose one.
\cref{thm:p15-bayes-reweighting,def:p17-finite-maxmin}. \\
A nominal interval comes from a poor sampling model & The model's coverage
calculation need not describe the actual procedure. & Check design,
dependence and distributions; use an exact or robust procedure with its
declared model class. No distribution-free guarantee follows from the word
robust. \cref{sec:p15-robust-hypothesis-testing}. \\
Some relevant causal variables are latent & Ordinary independent observed
kernel multiplication need not be valid. & Retain the latent model or use
the stated latent-graph identification machinery. Preserve query, support
and graph assumptions.
\cref{def:p16-causal-latent-projection,sec:p16-graph-uniform-identification}. \\
No effect is identified in the specified class & Compatible observational
models can give different effects. & Add justified structural restrictions,
collect appropriate intervention information, or report the unresolved
target. A more accurate observational estimate alone supplies no omitted
identification premise.
\cref{thm:p16-causal-nonidentification}. \\
An approximation lacks an error bound & The computed value does not yet
have the claimed accuracy. & Establish a remainder or numerical-error
bound on the actual domain. Relative error at a zero target needs a different
interpretation.
\cref{sec:p15-approximate-quantitative-inference,
sec:p12-numerical-error-and-inference}. \\
A response score is treated as a unique process measure & Several cognitive
accounts can produce the same coded response. & Supply process-sensitive
observations or a design distinguishing the accounts, with the task changes
recorded. \cref{def:p19-empirical-question,def:p19-evidence-record}. \\
A heuristic finds no candidate within a budget & No general nonexistence
conclusion follows without exhaustive-search conditions. & Use a complete
decision method for an eligible finite input, or retain a bounded search
conclusion. \cref{sec:p14-finder-proving,sec:p14-finder-deciding}. \\
\bottomrule
\end{longtable}

To compare two proposed methods, write the target in identical terms, check
that each method receives the same admissible information, and retain its
own output interpretation. A posterior risk, a significance level and a
robust preference can share numerical scales without sharing a guarantee.
If the methods assume different observation, causal or preference models,
the resulting disagreement may concern those models rather than a faulty
calculation. The linked definitions make the disputed premise inspectable.

For an empirical comparison, also fix task version, participants, scoring,
available information and the comparison design. The procedure in
\cref{def:p19-evidence-record} connects these to the reported conclusion;
it neither converts a convenience sample into a population estimate nor
turns an observational association into an intervention effect. A handbook
comparison remains useful when it identifies precisely which supplied
assumption a stronger conclusion would require.
